\section{Introduction}\label{sec:introduction}

Let $\theta\in(0,1)$ be irrational.  In the tracial representation let
$R_\theta$ be the von Neumann algebra generated by unitaries $U,V$ with
\[
 VU=e^{2\pi i\theta}UV,\qquad
 \tau(U^mV^n)=0\quad((m,n)\in\Z^2\setminus\{(0,0)\}),\qquad\tau(1)=1.
\]
This is the hyperfinite factor of type $\mathrm{II}_1$; the concrete
representation used below is the group-measure-space representation of
irrational rotation.  For a continuous function $f$ on the unit circle
$\T$, write $T_f=Uf(V)$.  A closed invariant subspace is
\emph{affiliated with $R_\theta$} when its orthogonal projection belongs
to $R_\theta$.  A nontrivial such subspace is therefore equivalent to
a projection $p\in R_\theta$ satisfying
\[
 0<\tau(p)<1,\qquad (1-p)T_fp=0.
\]
There is no requirement that $p$ commute with $T_f$.

Let $m$ denote normalized Haar measure on $\T$.  The
Fuglede--Kadison determinant of the scalar weight is
\[
 \Delta(f(V))=\exp\left(\int_\T\log|f(z)|\,dm(z)\right),
\]
with $\log0=-\infty$ and $\exp(-\infty)=0$.
Zhu, Fang, and Shi asked whether $T_f$ must have a nontrivial affiliated
invariant subspace if $f$ is nonzero almost everywhere, $|f|$ is
nonconstant, and this determinant is zero
\citep[p.~264, Question after Proposition~5.7]{zfs}.
We answer their question negatively for every prescribed irrational angle.

\begin{theorem}\label{thm:main}
For every irrational $\theta\in(0,1)$ there exists a continuous function
$f:\T\to[0,1]$ such that
\[
 f^{-1}(\{0\})=\{1\},\qquad \int_\T\log f\,dm=-\infty,
\]
and every projection $p\in R_\theta$ satisfying $(1-p)Uf(V)p=0$
is either $0$ or $1$.
\end{theorem}

The angle is fixed before the weight is constructed; no Diophantine
condition is imposed.  The conclusion concerns projections in the
specified containing factor.  Arbitrary invariant subspaces of its
Hilbert-space representation need not have their projections in that
factor.  Corollary~\ref{cor:quasinilpotent} also shows directly that the
constructed operator is nonzero and norm-quasinilpotent.

\subsection{Hyperinvariant subspaces and commutants}

For a bounded operator $T$ on a complex Hilbert space $\mathcal H$, write
$\{T\}'=\{A\in B(\mathcal H):AT=TA\}$.  A closed subspace is
\emph{hyperinvariant} for $T$ if every member of $\{T\}'$ leaves it
invariant.  An operator algebra is \emph{transitive} if its only common
closed invariant subspaces are $0$ and $\mathcal H$.

\begin{lemma}[Commutant algebra]\label{lem:commutant-algebra}
If $T\in B(\mathcal H)$ is nonscalar, then $\{T\}'$ is a proper,
strong-operator closed, unital complex algebra.  It is transitive exactly
when $T$ has no nontrivial hyperinvariant subspace.
\end{lemma}

\begin{proof}
The commutation equation is preserved by addition, multiplication and
scalar multiplication, and the identity commutes with $T$.  If
$A_\lambda\to A$ strongly and $A_\lambda T=TA_\lambda$, then for every
$\zeta\in\mathcal H$ both sides converge, giving $AT\zeta=TA\zeta$.
Thus the algebra is strongly closed.  If every bounded operator commuted
with $T$, then every rank-one orthogonal projection would commute with
$T$, so every nonzero vector would be an eigenvector.  Applying this to
$x,y,x+y$ for linearly independent vectors shows that the eigenvalues
are equal; hence $T$ would be scalar.  Properness follows.  The last
assertion is precisely the two definitions.
\end{proof}

The bridge to the containing von Neumann algebra is the elementary
observation recorded in \citet[Proposition~2.1]{dykema2005}.

\begin{lemma}\label{lem:hyperinvariant-projection}
Let $\mathcal M\subset B(\mathcal H)$ be a von Neumann algebra and
$T\in\mathcal M$.  The orthogonal projection onto any hyperinvariant
subspace for $T$ belongs to $\mathcal M$ and is invariant for $T$.
\end{lemma}

\begin{proof}
Every unitary $W\in\mathcal M'$ and its inverse commute with $T$.
Hyperinvariance therefore makes the subspace reducing for $W$.  Its
orthogonal projection commutes with all such unitaries and hence belongs
to $\mathcal M''=\mathcal M$.  Since $T\in\{T\}'$, its range is
invariant under $T$.
\end{proof}

\begin{corollary}\label{cor:transitive-commutant}
Every separable infinite-dimensional complex Hilbert space carries a
nonzero norm-quasinilpotent operator with no nontrivial hyperinvariant
subspace.  Its commutant is a proper strongly closed transitive unital
algebra.
\end{corollary}

\begin{proof}
The tracial representation of $R_\theta$ is separable and
infinite-dimensional.  Apply Theorem~\ref{thm:main},
Corollary~\ref{cor:quasinilpotent}, and
Lemma~\ref{lem:hyperinvariant-projection} there.  A nonzero
quasinilpotent operator is nonscalar, so
Lemma~\ref{lem:commutant-algebra} applies.  Unitary equivalence transfers
these properties to any separable infinite-dimensional complex Hilbert
space.
\end{proof}

These conclusions do not exclude ordinary invariant subspaces in the
representation.  A separate companion gives a direct $2$-adic weighted
shift with a transitive commutant, including an explicit norm estimate,
backward intertwiners and a description of its polar conjugates
\citep{direct-companion}.  It uses only
Lemmas~\ref{lem:commutant-algebra} and~\ref{lem:scalar-multiplier-fields}
from this paper.  The present proof uses no result from that companion.

\begin{remark}\label{rem:ckl}
Corollary~\ref{cor:transitive-commutant} contradicts the hyperinvariant-subspace assertion of
\citet[Corollary~2.8, p.~801]{cho-ko-lee} at $\mu=1$: for the nonzero
quasinilpotent $T$ just constructed and any conjugation $C$, their stated
hypotheses apply to $I+T$, but $\{I+T\}'=\{T\}'$, so $I+T$ has no
nontrivial hyperinvariant subspace.  Only that hyperinvariant-subspace assertion is at issue here.
\end{remark}

\subsection{Prior work}

The determinant of Fuglede and Kadison \citep{fuglede-kadison} and
Brown's spectral distribution \citep{brown} form a major route to
invariant subspaces in finite factors.  Haagerup and Schultz associate an
affiliated hyperinvariant projection to each Borel subset of the plane,
with trace equal to its Brown mass \citep[Theorem~1.1]{haagerup-schultz}.
Thus a Brown measure that is not a point mass gives a nontrivial
projection.  For a point-mass Brown measure, the trace formula gives
only trivial projections.

Point-mass Brown measure can nevertheless coexist with many invariant
subspaces.  Dykema and Haagerup construct a one-parameter family of
nontrivial affiliated hyperinvariant subspaces for the quasinilpotent
DT-operator \citep[Theorems~2.2 and~2.9]{dykema-haagerup}.
Thus degeneration of Brown measure to a point mass does not itself
exclude invariant projections.
Tucci also constructed quasinilpotent generators of the hyperfinite factor
with nontrivial affiliated invariant projections
\citep[Theorems~3.1 and~3.3 and the discussion following Theorem~3.3]{Tucci2008}.
These examples further distinguish quasinilpotence from the absence of
factor-invariant projections.

For weighted irrational rotations, the measurable-action theorem of
Dykema and Schultz identifies the Brown measure as uniform measure on
the circle whose radius is the determinant
\citep[Theorem~5.4]{dykema-schultz}.  Zhu, Fang, and Shi give a direct
proof, explicitly crediting that earlier result
\citep[Theorem~5.3]{zfs}.  At determinant zero this measure is the point
mass at zero, so the Brown-measure projections are trivial.

Continuous weights with one zero and determinant zero were already
exhibited by Zhu, Fang, and Shi \citep[pp.~262--263]{zfs}.
The additional conclusion of Theorem~\ref{thm:main} is the exclusion of
every nontrivial invariant projection for a suitable weight at each
prescribed angle.  Zhu, Fang, and Shi's spectral-radius formula
\citep[Theorem~2.3]{zfs} implies quasinilpotence in this case.
We give an elementary proof of this special case, using only uniform
averaging of continuous functions under irrational rotation.
Their positive-determinant invariant-subspace result
\citep[Corollary~5.5]{zfs} has a hypothesis complementary to the one
studied here.

\subsection{Proof strategy}

The obstruction begins with a measurable weight.  Suppose its weighted
rotation had a nontrivial invariant projection.  In the shift
representation, this projection becomes a field of projections.  The
finite trace and ergodicity force its central diagonal entry to lie
strictly between zero and one almost everywhere.  Thus both the
projection and its complement have vectors with a nonzero central
coordinate.

Invariance imposes weighted orthogonality identities on these vectors.
We first obtain the identities along rotation orbits and then pass them
to suitably coupled pairs of points.  The weight has drops that are rare
enough to leave a chosen finite window under control, but strong enough
to suppress contributions outside that window.  The central coordinates
force at least one finite convolution coefficient to stay away from
zero.  The suppressed tail cannot cancel it, contradicting the exact
orthogonality identity.  Conditioning separately on each endpoint
controls that tail even though it has infinitely many terms.

The measurable weight is then replaced by a continuous weight with one
zero.  Smoothing and moving the logarithmic drops changes their sum by
a constant plus a measurable coboundary.  This relation supplies an
intertwining multiplier, possibly unbounded.  A domain and trace argument
shows that an invariant projection for the continuous operator would
produce one for the measurable operator with exactly the same trace.
This transfers the obstruction to the desired continuous weight.

The construction realizing this scheme is as follows.
Use the additive circle $X=\R/\Z$, with Haar probability $\mu$, and write
$\sigma x=x+\theta$.  Write $\{t\}\in[0,1)$ for fractional part.
We construct divisible integers
$1=q_0\mid q_1\mid q_2\mid\cdots$ and positive costs $c_i$.
The measurable model is the weight
\[
 g_*(x)=\sum_{i\ge1}c_i\lfloor\{q_ix\}+\{q_i\theta\}\rfloor,
 \qquad f_*(x)=e^{-g_*(x)}.
\]
The series is finite almost everywhere although its integral diverges.
Large divisible frequencies allow its trial phases to be realized at the
fixed angle $\theta$, while finitely many sampling tests approximate
independent Haar coordinates.  The circle characters themselves are
dependent.  A mixed-radix expansion supplies independent digits, from
which we build the ordered couplings with controlled marginal densities.
The prescribed-marginal viewpoint belongs to the classical coupling
theory of Strassen and Edwards \citep{Strassen1965,Edwards1978}; here
we also need approximation by positive orbit pairs.

Section~\ref{sec:parameters} fixes the frequencies and sampling powers.
Sections~\ref{sec:fields} and~\ref{sec:couplings} obtain the projection
identities on the ordered pairs.  Section~\ref{sec:obstruction} proves
and applies the convolution obstruction, and
Section~\ref{sec:continuity} constructs the continuous weight and proves
the trace-preserving transfer.

Appendix~\ref{p:sec:product} gives a complete alternative realization on
an infinite product of circles.  Independent coordinates yield a
finite-coordinate orbit-sampling theorem and a second application of the
same convolution obstruction.  That model also separates the atomic
algebras generated by finitely many polar conjugates from the diffuse
algebra generated by the full orbit.  The prescribed-angle circle proof
is complete before the appendix.



All identities involving measurable fields hold on appropriate conull
sets.  When integer iterates are used, we intersect the countably many
translates of these sets.  The constants implicit in $O(\cdot)$ below
are absolute unless a fixed construction stage is specified.  Every new
ingredient needed for the proof is established in the paper.
